% ============================================================
%  ACADEMIC COVER LETTER TEMPLATE
%  Suitable for: Academic positions, fellowships, graduate school
%
%  LaTeX concepts covered:
%    \documentclass{letter}  --- LaTeX's built-in letter class
%    \address{}              --- your return address
%    \signature{}            --- your sign-off name
%    \begin{letter}{addr}    --- recipient address
%    \opening{}              --- salutation (Dear ...)
%    \closing{}              --- closing phrase (Sincerely, ...)
%    \today                  --- inserts today's date automatically
%    \encl{}                 --- list of enclosed documents
%    \ps{}                   --- postscript
%
%  The letter class handles all spacing and layout conventions
%  automatically. You focus entirely on the content.
%
%  Instructions:
%    1. Replace ALL placeholder text with your own information.
%    2. Keep each paragraph focused: one idea per paragraph.
%    3. Click Recompile (or Ctrl+Enter) to preview the PDF.
% ============================================================

% ------ DOCUMENT CLASS ------
% The "letter" class is specifically designed for formal correspondence.
% It is different from "article" — it provides letter-specific commands.
\documentclass[12pt]{letter}

\usepackage[T1]{fontenc}
\usepackage[utf8]{inputenc}
\usepackage[top=2.5cm, bottom=2cm, left=3cm, right=3cm]{geometry}
\usepackage[hidelinks]{hyperref}

% \longindentation controls how far the closing/signature is indented.
% Setting it to 0pt aligns the closing to the left margin.
\longindentation=0pt


% ------ YOUR RETURN ADDRESS ------
% \address sets the sender's address printed at the top of the letter.
% Use \\ to separate lines.
%
% IMPORTANT: if the next line starts with [ LaTeX misreads it as a spacing
% argument. Use \\{} to block that — the empty {} acts as a separator.
\address{
  Your Full Name \\{}
  Department of Your Department \\{}
  University Name \\{}
  City, Province, Postal Code \\
  \href{mailto:your.email@university.edu}{your.email@university.edu}
}

% ------ YOUR SIGNATURE ------
% \signature sets the name that appears after the closing phrase.
% Same rule applies inside \signature: use \\{} before any line starting with [
\signature{Your Full Name  \\{} PhD Candidate in Your Field}


% ============================================================
%  DOCUMENT BEGINS HERE
% ============================================================
\begin{document}

% ------ BEGIN THE LETTER ------
% The argument to \begin{letter}{...} is the recipient's address.
\begin{letter}{
  Recipient Name \\{}
  Title, e.g., Chair of the Search Committee \\{}
  Department of Department Name \\{}
  University Name \\{}
  City, Country
}

% \today inserts the current date automatically.
% You can replace it with a fixed date, e.g., April 14, 2025.
\date{\today}

% ------ SALUTATION ------
% \opening{} produces the "Dear ..." line.
\opening{Dear Professor Surname,}


% ------ PARAGRAPH 1: INTRODUCTION ------
% State who you are, what position you are applying for, and how you found it.
I am writing to apply for the position of Postdoctoral Researcher in Your Field
at University Name, as advertised on Website or Journal Name. I am currently
completing my PhD in Your Field at Your University under the supervision of
Professor Supervisor Name, and I expect to defend my dissertation in Month Year.


% ------ PARAGRAPH 2: RESEARCH OVERVIEW ------
% Summarise your research area and its significance in 3--4 sentences.
My dissertation investigates describe your research topic in plain language.
This work addresses describe the problem or gap in the field by describing
your approach or method. My findings demonstrate summarise a key result or
contribution. This research has resulted in X publications and presentations
at conferences including Conference Name.


% ------ PARAGRAPH 3: FIT WITH THE DEPARTMENT ------
% Explain specifically why you want this position and how you fit.
I am drawn to University Name because of a specific reason such as a faculty
member's work, a lab, or a department strength. My research on your topic
complements the work of Professor Name on their topic, and I am excited by
the possibility of describe a potential collaboration or shared direction.


% ------ PARAGRAPH 4: TEACHING ------
% Briefly mention your teaching experience and philosophy.
I have teaching experience from my role as a Teaching Assistant for Course Name
at University Name, where I briefly describe your responsibility. I am prepared
to teach courses in Area One and Area Two, and I am eager to develop new courses
in an emerging area relevant to the department.


% ------ PARAGRAPH 5: CLOSING ------
I have enclosed my curriculum vitae, research statement, teaching statement,
and the names and contact details of three referees. I welcome the opportunity
to discuss my work with you. Thank you sincerely for your time and consideration.


% ------ CLOSING ------
% \closing{} produces the sign-off phrase. \signature{} text appears below it.
\closing{Sincerely,}


% ------ ENCLOSURES ------
% \encl{} lists the documents enclosed with the letter.
\encl{Curriculum Vitae \\
      Research Statement \\
      Teaching Statement \\
      Contact information for three referees}

% ------ POSTSCRIPT (optional) ------
% \ps{} adds a postscript after the signature. Remove if not needed.
% \ps{Please feel free to contact my primary referee, Prof.\ Name, at email.}

\end{letter}

% ============================================================
%  DOCUMENT ENDS HERE
% ============================================================
\end{document}
